\section*{الفصل \ref{ch:b1:multivar} --- دوال المتغيرين}

\begin{solution}{exo:b1:multivar:1}
$\dfrac{\partial f}{\partial x} = 2xy + y\,\eu^{xy}$،
$\dfrac{\partial f}{\partial y} = x^2 + x\,\eu^{xy}$.

$\dfrac{\partial g}{\partial x} = \dfrac{2x}{x^2+y^2}$،
$\dfrac{\partial g}{\partial y} = \dfrac{2y}{x^2+y^2}$.

$\dfrac{\partial h}{\partial x} = \dfrac{-y/x^2}{1 + y^2/x^2} =
\dfrac{-y}{x^2+y^2}$،
$\dfrac{\partial h}{\partial y} = \dfrac{x}{x^2+y^2}$.
\end{solution}

\begin{solution}{exo:b1:multivar:2}
\omterm{prop:b1:vspaces:coordinates}{بالإحداثيات} القطبية ($\rho \to 0$):

$\abs{f} = \dfrac{\rho^3\abs{\cos^2\theta\sin\theta}}{\rho^2} \leq
\rho \to 0$: \omterm{def:b1:continuity:continuous}{متصلة}.

$g = \cos\theta\sin\theta$: مستقل عن $\rho$، ويأخذ قيمًا
مختلفة على امتداد أنصاف أقطار مختلفة (قارن
\cref{ex:b1:multivar:trap}): فلا نهاية، وهي غير \omterm{def:b1:continuity:continuous}{متصلة}.

$\abs h \leq \dfrac{\rho^3(\abs{\cos^3\theta} +
\abs{\sin^3\theta})}{\rho^2} \leq 2\rho \to 0$: \omterm{def:b1:continuity:continuous}{متصلة}.
\end{solution}

\begin{solution}{exo:b1:multivar:3}
على امتداد $y = mx$: $f(x, mx) = \dfrac{m x^3}{x^4 + m^2 x^2} =
\dfrac{mx}{x^2 + m^2} \to 0$ (من أجل $m \neq 0$؛ وعلى امتداد $y = 0$ و
محور $y$، $f = 0$). ومنه فكل نهاية على مستقيم هي $0$. لكن على
القطع المكافئ $y = x^2$:
\[
f(x, x^2) = \frac{x^4}{x^4 + x^4} = \frac12 :
\]
وللمتتالية $\bigl(\frac1n, \frac{1}{n^2}\bigr) \to (0,0)$ لدينا
$f \to \frac12 \neq 0$. فهي غير \omterm{def:b1:continuity:continuous}{متصلة}: إذ لا تكفي المستقيمات
لاختبار نهايات المتغيرين.
\end{solution}

\begin{solution}{exo:b1:multivar:4}
$\frac{\partial f}{\partial x} = 3x^2y^2 + y\cos(xy)$؛ ثم
\[
\frac{\partial^2 f}{\partial y\,\partial x}
= 6x^2 y + \cos(xy) - xy\sin(xy) .
\]
$\frac{\partial f}{\partial y} = 2x^3 y + x\cos(xy)$؛ ثم
\[
\frac{\partial^2 f}{\partial x\,\partial y}
= 6x^2 y + \cos(xy) - xy\sin(xy) :
\]
فهما متساويان، كما يعد شوارتز.
\end{solution}

\begin{solution}{exo:b1:multivar:5}
حسب \cref{thm:b1:multivar:chain} مع $(x, y) = (\cos t, \sin t)$:
\[
g'(t) = -\sin t\,\frac{\partial f}{\partial x}(\cos t, \sin t)
+ \cos t\,\frac{\partial f}{\partial y}(\cos t, \sin t)
= \bigl\langle \nabla f,\ (-\sin t, \cos t)\bigr\rangle .
\]
وعند قيمة حدّية للمقدار $g$، $g'(t) = 0$: ومنه \omterm{def:b1:multivar:partial}{فالتدرج} $\nabla f$ متعامد مع
متجهة مماس الدائرة $(-\sin t, \cos t)$، ومنه موازٍ
لمتجهة نصف القطر $(\cos t, \sin t)$ (لأن \omterm{def:b1:euclid:orthogonal}{المتمم المتعامد}
لمتجهة واحدية في المستوي هو المستقيم الذي تولّده، مأخوذًا
عموديًا). وهذا أبسط مثال على مضاعف
لاغرانج.
\end{solution}

\begin{solution}{exo:b1:multivar:6}
$f$: $\nabla f = (2x + y - 3,\; x + 2y) = 0$: أي $y = -\frac x2$ و
$2x - \frac x2 = 3$: ومنه $x = 2$ و $y = -1$. والرتبة الثانية: $r = 2$ و $s =
1$ و $t = 2$: ومنه $rt - s^2 = 3 > 0$ و $r > 0$: أي قيمة صغرى موضعية (بل شاملة ---
لأنها تربيعية) عند $(2, -1)$، وقيمتها $f(2,-1) = -3$.

$g$: $\nabla g = (2x,\; -2y + 4) = 0$: أي النقطة $(0, 2)$؛ و $r = 2$ و $s
= 0$ و $t = -2$: ومنه $rt - s^2 = -4 < 0$: أي سرج.

$h$: $\nabla h = \bigl(4x^3 - 4(x - y),\; 4y^3 + 4(x-y)\bigr) = 0$.
وبالجمع: $x^3 + y^3 = 0$، ومنه $y = -x$؛ وبالتعويض: $4x^3 - 8x = 0$:
$x \in \{0, \pm\sqrt2\}$. فالنقاط الحرجة: $(0,0)$ و
$(\sqrt2, -\sqrt2)$ و $(-\sqrt2, \sqrt2)$.
وعند $(\pm\sqrt2, \mp\sqrt2)$: $r = 12x^2 - 4 = 20$ و $s = 4$ و $t =
20$: ومنه $rt - s^2 > 0$ و $r > 0$: أي قيمتان صغريان موضعيتان (بالقيمة $h = 4 + 4 - 16 =
-8$). وعند $(0,0)$: $r = t = -4$ و $s = 4$: ومنه $rt - s^2 = 0$: فهو صامت.
وافحص: $h(x, x) = 2x^4 > 0$ و $h(x, -x) = 2x^4 - 8x^2 < 0$ من أجل
$x \neq 0$ صغير: أي الإشارتان في كل \omterm{def:b1:topology:open}{جوار} --- فلا قيمة حدّية
عند المبدأ.
\end{solution}

\begin{solution}{exo:b1:multivar:7}
$x^2 - y^2$: المشتقتان الجزئيتان الثانيتان $2$ و $-2$: فالمجموع $0$. و $xy$: تنعدم
المشتقتان الثانيتان المحضتان كلتاهما. و $\eu^x\cos y$: $\frac{\partial^2}{\partial
x^2} = \eu^x\cos y$ و $\frac{\partial^2}{\partial y^2} =
-\eu^x\cos y$: فالمجموع $0$. و $\ln(x^2+y^2)$: من
\cref{exo:b1:multivar:1}،
\[
\frac{\partial^2}{\partial x^2}\ln(x^2+y^2)
= \frac{2(x^2+y^2) - 4x^2}{(x^2+y^2)^2}
= \frac{2(y^2 - x^2)}{(x^2+y^2)^2},
\]
ونظير $y$ هو مقابله: فالمجموع $0$.
\end{solution}

\begin{solution}{exo:b1:multivar:8}
\emph{\omterm{def:b1:linmaps:projection}{بالإسقاط}:} للمستوي $P$ ناظمٌ $n = (1, 2, -1)$ وهو
يمرّ بالمقدار $A = (4, 0, 0)$. وأقرب نقطة إلى المبدأ هي
$O$ مسقطًا: $p = O + \frac{\langle A - O, n\rangle}{\norm
n^2}\,n = \frac{4}{6}(1,2,-1) = \bigl(\frac23, \frac43,
-\frac23\bigr)$، على مسافة $\frac{4}{\sqrt 6}$. (والصيغة: المسافة
من المبدأ إلى المستوي $\langle X, n \rangle = c$ هي
$\frac{\abs c}{\norm n}$ مع $c = 4$.)

\emph{وبالتصغير:} $f(x,y) = x^2 + y^2 + (x + 2y - 4)^2$. وبكتابة
$w = x + 2y - 4$ من أجل العامل الأخير:
\[
\nabla f = \bigl(2x + 2w,\; 2y + 4w\bigr) = 0
\iff x = -w \text{ و } y = -2w .
\]
وبالتعويض في تعريف $w$: $w = -w - 4w - 4$، ومنه $w =
-\frac23$، فينتج $x = \frac23$ و $y = \frac43$ و $z = w =
-\frac23$. أي النقطة نفسها كما \omterm{def:b1:linmaps:projection}{بالإسقاط}، على مسافة
$\sqrt{\frac49 + \frac{16}{9} + \frac49} = \frac{2\sqrt 6}{3} =
\frac{4}{\sqrt6}$؛ وهي قيمة صغرى لأن $f \to \infty$ عند اللانهاية
(لأنها تربيعية معيَّنة موجبة زائد حدود خطية).
\end{solution}

\begin{solution}{exo:b1:multivar:9}
مع $\rho = \sqrt{x^2+y^2}$: $\frac{\partial \rho}{\partial x} =
\frac{x}{\rho}$، ومنه
\[
\frac{\partial f}{\partial x} = \varphi'(\rho)\,\frac{x}{\rho},
\qquad
\frac{\partial^2 f}{\partial x^2}
= \varphi''(\rho)\,\frac{x^2}{\rho^2}
+ \varphi'(\rho)\,\Bigl(\frac{1}{\rho} -
\frac{x^2}{\rho^3}\Bigr),
\]
(بقاعدتَي القسمة والسلسلة). وبإضافة \omterm{def:b1:logic:statement}{عبارة} $y$ المتناظرة:
تتجمّع حدود $\varphi''$ لتعطي $\frac{x^2 + y^2}{\rho^2} = 1$، و
تتجمّع حدود $\varphi'$ لتعطي $\frac{2}{\rho} - \frac{x^2 +
y^2}{\rho^3} = \frac{1}{\rho}$:
\[
\Delta f = \varphi'' + \frac{\varphi'}{\rho} .
\]
والتوافقية الشعاعية: حُلَّ $\varphi'' + \frac{\varphi'}{\rho} = 0$:
فالمعادلة من الرتبة الأولى $u' + \frac u\rho = 0$ من أجل $u = \varphi'$
تعطي $u = \frac{c}{\rho}$ (\cref{thm:b1:diffeq:homogeneous1})،
ثم $\varphi = c\ln\rho + d$. ومنه فالدوال التوافقية الشعاعية هي
$c\,\ln\sqrt{x^2 + y^2} + d$ --- أي الكمون اللوغاريتمي في
\cref{exo:b1:multivar:7} والثوابت، لا غير.
\end{solution}

\begin{solution}{exo:b1:multivar:10}
من أجل $z = xy$ عند $(1,1)$: المشتقتان الجزئيتان $y = 1$ و $x = 1$، فالمستوي
المماس $z = 1 + (x - 1) + (y - 1) = x + y - 1$. ويعني المستوي المماس
الأفقي انعدام المشتقتين الجزئيتين معًا، أي $\nabla f = 0$:
ومنه، حسب \cref{ex:b1:multivar:extrema}، تكون تلك بالضبط النقاط الحرجة
$(0, 0)$ و $(1, 1)$، بمستويين أفقيين $z = 0$ و $z =
-1$. \omterm{def:b1:logic:statement}{فعبارة} «المستوي المماس الأفقي» و«\omterm{thm:b1:multivar:critical}{النقطة الحرجة}» هما
المفهوم نفسه، منظورًا إليه على المنحني وفي الصيغة.
\end{solution}

\begin{solution}{exo:b1:multivar:11}
\begin{enumerate}
  \item من أجل $y$ مثبَّت، تكون للدالة في متغير واحد $x \mapsto
        f(x,y)$ مشتقةٌ معدومة على $\R$، ومنه فهي ثابتة
        (\cref{thm:b1:derivative:mvt}): أي $f(x, y) = f(0, y)$ من أجل
        كل $x$: ومنه فالدالة $f$ لا تتعلق إلا بالمقدار $y$.
  \item لتكن $g(u, v) = f\bigl(\frac{u+v}2, \frac{u-v}2\bigr)$.
        وبقاعدة السلسلة (\cref{thm:b1:multivar:chain}، مطبَّقةً
        في المتغير $v$ مع تجميد $u$):
        \[
        \frac{\partial g}{\partial v}
        = \frac12\,\frac{\partial f}{\partial x}
        - \frac12\,\frac{\partial f}{\partial y} = 0 .
        \]
        وحسب (1)، لا يتعلق $g$ إلا بالمقدار $u$: أي $g(u, v) = \varphi(u)$
        مع $\varphi$ من الصنف $C^1$، وبعكس تغيير
        المتغيرات ($u = x + y$ و $v = x - y$)،
        \[
        f(x, y) = \varphi(x + y).
        \]
        وبالعكس تحقق كل $f$ كهذه $f_x = f_y =
        \varphi'$: ومنه فالحلول هي بالضبط الدوال من الصنف $C^1$
        في $x + y$.
\end{enumerate}
\end{solution}

\begin{solution}{exo:b1:multivar:12}
على القرص المفتوح، ستكون القيمة الحدّية حرجة: و $\nabla f = (y,
x) = 0$ عند $(0,0)$ فقط، حيث $f = 0$؛ وهي سرج ($f(\pm
\varepsilon, \pm\varepsilon) = \varepsilon^2 > 0 > -\varepsilon^2
= f(\pm\varepsilon, \mp\varepsilon)$)، ومنه فلا قيمة حدّية هناك. وحسب
مبرهنة القيم الحدّية المقبولة تُبلغ الحواصر،
وبالضرورة على الدائرة الحدّية. وهناك، مع
\cref{exo:b1:multivar:5}،
\[
f(\cos t, \sin t) = \cos t\sin t = \frac{\sin 2t}{2}
\in \intcc{-\tfrac12}{\tfrac12},
\]
بقيمة عظمى $\frac12$ عند $t = \frac\pi4, \frac{5\pi}4$ (أي النقطتان
$\pm\frac{1}{\sqrt2}(1,1)$) وقيمة صغرى $-\frac12$ عند $t =
\frac{3\pi}4, \frac{7\pi}4$ (أي النقطتان $\pm\frac{1}{\sqrt2}(1,-1)$).
فالقيمة العظمى الشاملة $\frac12$ والصغرى الشاملة $-\frac12$.
\end{solution}

\begin{solution}{pb:b1:multivar:1}
\textbf{1.} بنشر الطرف الأيمن:
\[
r\Bigl(h + \frac srk\Bigr)^{\!2} + \frac{rt - s^2}{r}k^2
= rh^2 + 2shk + \frac{s^2}{r}k^2 + \frac{rt - s^2}{r}k^2
= q(h,k) .
\]
فإذا كان $rt - s^2 > 0$: حمل المربعان كلاهما إشارة عامل $r$،
ويفرض $q(h,k) = 0$ أن $k = 0$ ثم $h = 0$: أي إشارة $r$ القطعية
خارج المبدأ. وإذا كان $rt - s^2 < 0$: فإن $q(1, 0) = r$ بينما
$q\bigl(-\frac sr, 1\bigr) = \frac{rt - s^2}{r}$ له
الإشارة المقابلة: فتقع الإشارتان.

\textbf{2.} إذا كان $r = 0 \neq t$، فبادل بين دورَي $h$ و $k$
(بإتمام المربع في $k$)، بالنتائج نفسها؛ ولاحظ
$rt - s^2 = -s^2 < 0$ ما إن يكن $s \neq 0$، وفعلًا يأخذ $q =
2shk + tk^2$ عندئذ الإشارتين ($k$ صغيرًا إزاء $h$). وإذا كان
$r = t = 0$: $q = 2shk$، فالإشارتان إذا وفقط إذا كان $s \neq 0$، أي إذا وفقط إذا كان
$rt - s^2 = -s^2 < 0$. وتلخيصًا: يأخذ $q$ الإشارتين $\iff rt -
s^2 < 0$؛ وينعدم $q$ عند المبدأ فقط (معيَّنًا) $\iff rt -
s^2 > 0$ (وهذا يفرض $r \neq 0$، لأن $r = 0$ يعطي $q(1,0) =
0$).

\textbf{3.} مع $f_x = rx + sy + \beta$ و $f_y = sx + ty +
\gamma$، يعطي النشر المباشر
\[
f(a + h, b + k) = f(a, b) + h\,f_x(a,b) + k\,f_y(a,b)
+ \tfrac12 q(h, k),
\]
دون أيّ حدود أعلى (لأن الدالة \omterm{def:b1:poly:def}{كثير حدود} من الدرجة
$2$). وعند \omterm{thm:b1:multivar:critical}{نقطة حرجة} ينعدم الجزء الخطي: أي $f(X_0 + H)
- f(X_0) = \frac12 q(H)$ \emph{بالضبط}، ومنه تصنّف دراسة الإشارة في
السؤالين 1--2: قيمة صغرى شاملة قطعية ($rt - s^2 >
0$ و $r > 0$)، وقيمة عظمى شاملة قطعية ($rt - s^2 > 0$ و $r < 0$)،
وسرج --- بالإشارتين في كل \omterm{def:b1:topology:open}{جوار} --- عندما $rt - s^2 <
0$. وهذا يبرهن على \cref{met:b1:multivar:monge} من أجل الدوال
التربيعية.

\textbf{4.} للشكل $q - c(h^2 + k^2)$ المعطيات $r - c$ و $s$ و
$t - c$. ومع $c = \frac{rt - s^2}{r + t}$ (ولاحظ $t > 0$ لأن
$rt > s^2 \geq 0$ و $r > 0$):
\[
(r - c)(t - c) - s^2 = rt - s^2 - c(r + t) + c^2 = c^2 \geq 0,
\]
و $r - c \geq 0$ ($c \leq r \iff rt - s^2 \leq r^2 + rt$،
وهو صحيح). فإذا كان $r - c > 0$، بيّن إتمام المربع في السؤال 1 أن
$q - c(h^2 + k^2) \geq 0$؛ وإذا كان $r - c = 0$، فإن $s = 0$ (من
$-s^2 = -c^2 \leq \dots$ الذي يفرض أن يكون المقدار المعروض
$\geq 0$ وعامله الأول $0$) ويكون الشكل $(t - c)k^2 \geq
0$. وفي الحالتين $q(h,k) \geq c(h^2 + k^2)$.

\textbf{5.} حسب السؤالين 3--4، $f(X) = f(X_0) + \frac12 q(X -
X_0) \geq f(X_0) + \frac c2\norm{X - X_0}^2 \to +\infty$ عندما
$\norm X \to \infty$: ومنه فالدالة $f$ قاسرة، والمتراجحة
قطعية من أجل $X \neq X_0$: ومنه \omterm{thm:b1:multivar:critical}{فالنقطة الحرجة} هي المصغِّر
الشامل الوحيد. وأمّا من أجل التكعيبي $f = x^3 + y^3 - 3xy$ في
\cref{ex:b1:multivar:extrema}، فلا تصحّ نتيجة كهذه: $f(x,0)
= x^3 \to -\infty$، والقيمة الصغرى الموضعية عند $(1,1)$ ليست
شاملة --- فما فشل هو ضبط النشر التربيعي.

\textbf{6.} بنشر مربع المعيار:
\[
f(u,v) = u^2\norm{C_1}^2 + 2uv\,\langle C_1, C_2\rangle +
v^2\norm{C_2}^2 - 2u\,\langle C_1, b\rangle - 2v\,\langle C_2,
b\rangle + \norm b^2 ,
\]
أي دالة تربيعية في $(u, v)$ جزؤها التربيعي هو
$\frac12(ru^2 + 2suv + tv^2)$ مع $r = 2\norm{C_1}^2$ و $s =
2\langle C_1, C_2\rangle$ و $t = 2\norm{C_2}^2$.

\textbf{7.} $rt - s^2 = 4\bigl(\norm{C_1}^2\norm{C_2}^2 -
\langle C_1, C_2\rangle^2\bigr) \geq 0$ بكوشي--شوارتز
(\cref{thm:b1:euclid:cs})، مع المساواة بالضبط عندما يكون $C_1, C_2$
متناسبين (أو يكون أحدهما معدومًا)، أي عندما يكون الزوج
مرتبطًا. ومنه $rt - s^2 > 0 \iff (C_1, C_2)$ \omterm{def:b1:vspaces:free}{حرة}.

\textbf{8.} يُقرأ $\nabla f = 0$ هكذا
\[
\norm{C_1}^2 u + \langle C_1, C_2\rangle v = \langle C_1,
b\rangle,
\qquad
\langle C_1, C_2\rangle u + \norm{C_2}^2 v = \langle C_2,
b\rangle,
\]
أي \omterm{pb:b1:multivar:1}{المعادلات الناظمية} بمصفوفة غرام على اليسار؛ وهي تقول
$\langle uC_1 + vC_2 - b,\ C_i\rangle = 0$ من أجل $i = 1, 2$، أي
$b - p \perp \operatorname{Vect}(C_1, C_2)$ من أجل $p = uC_1 +
vC_2$. وحسب الجزء 1 (السؤالان 3 و 5) تكون \omterm{thm:b1:multivar:critical}{النقطة الحرجة} الوحيدة
هي المصغِّر الشامل الوحيد، وحسب
\cref{thm:b1:euclid:projection} يعيّن التمييز «$p \in
F$ و $b - p \perp F$» المقدارَ $p$ إسقاطًا متعامدًا
للمقدار $b$ على $F = \operatorname{Vect}(C_1, C_2)$.

\textbf{9.} $\norm{b - p}^2 = \norm b^2 - 2\langle b, p\rangle +
\norm p^2$، ويعطي $\langle p, b - p\rangle = 0$ أن $\norm p^2 =
\langle p, b\rangle$: ومنه فالقيمة الصغرى تساوي $\norm b^2 - \langle p,
b\rangle$. وبفيثاغورس: $\norm b^2 = \norm p^2 + \norm{b - p}^2$
--- فتنقسم متجهة المعطيات إلى جزئه المفسَّر $p$ وجزئه
الباقي $b - p$، وهما متعامدان.

\textbf{10.} $E(\alpha, \beta) = \norm{\alpha C_1 + \beta C_2 -
b}^2$ مع $C_1, C_2, b$ المذكورة؛ وتكون \omterm{pb:b1:multivar:1}{المعادلات الناظمية} في
السؤال 8، مركّبةً مركّبة،
\[
n\alpha + \Bigl(\sum x_i\Bigr)\beta = \sum y_i,
\qquad
\Bigl(\sum x_i\Bigr)\alpha + \Bigl(\sum x_i^2\Bigr)\beta = \sum
x_i y_i .
\]

\textbf{11.} بالقسمة على $n$: $\alpha + \beta\bar x = \bar y$
و $\alpha\bar x + \beta\bigl(v_x + \bar x^2\bigr) = c_{xy} +
\bar x\bar y$. وبتعويض $\alpha = \bar y - \beta\bar x$ في
الثانية: $\beta v_x = c_{xy}$، ومنه
\[
\beta = \frac{c_{xy}}{v_x}, \qquad \alpha = \bar y - \beta\bar
x ,
\]
وتقول المعادلة الأولى بالضبط إن $(\bar x, \bar y)$
تقع على المستقيم.

\textbf{12.} $v_x = \frac1n\sum(x_i - \bar x)^2 \geq 0$، وهو معدوم
إذا وفقط إذا ساوى كل $x_i$ المقدارَ $\bar x$. و $rt - s^2 = 4\bigl(n\sum
x_i^2 - (\sum x_i)^2\bigr) = 4n^2 v_x$: ومنه فشرط الحرية في السؤال 7
هو $v_x > 0$، أي ألّا تتساوى كل $x_i$.

\textbf{13.} مع $\alpha = \bar y - \beta\bar x$، مركّز
المعطيات ($\tilde x_i = x_i - \bar x$ و $\tilde y_i = y_i - \bar y$):
\[
E(\alpha, \beta) = \sum_i(\tilde y_i - \beta\tilde x_i)^2
= n\,v_y - 2\beta\,n\,c_{xy} + \beta^2 n\,v_x ,
\]
وهو أصغري عند $\beta = c_{xy}/v_x$ بالقيمة $E_{\min} =
n\bigl(v_y - \frac{c_{xy}^2}{v_x}\bigr) = n v_y(1 - \rho^2)$.
ولأن $E_{\min} \geq 0$ و $v_y > 0$: يكون $\rho^2 \leq 1$، و
$\abs\rho = 1$ إذا وفقط إذا كان $E_{\min} = 0$، أي إذا وفقط إذا انعدم كل
باقٍ: أي إن النقاط تقع على المستقيم بالضبط.

\textbf{14.} $n = 4$: $\bar x = \frac64 = 1.5$ و $\bar y = 2$ و
$\sum x_i^2 = 14$ ومنه $v_x = 3.5 - 2.25 = 1.25$، و $\sum x_iy_i =
0 + 1 + 6 + 12 = 19$ ومنه $c_{xy} = 4.75 - 3 = 1.75$. ومنه
\[
\beta = \frac{1.75}{1.25} = 1.4,
\qquad
\alpha = 2 - 1.4\times1.5 = -0.1 :
\qquad y = 1.4\,x - 0.1 .
\]
والقيم المقايَسة $-0.1,\ 1.3,\ 2.7,\ 4.1$؛ والبواقي $0.1,\ -0.3,\
0.3,\ -0.1$؛ و $E_{\min} = 0.01 + 0.09 + 0.09 + 0.01 = 0.2$
(وللتحقق: $v_y = \frac{26}4 - 4 = 2.5$ و $4(2.5 -
\frac{1.75^2}{1.25}) = 4(2.5 - 2.45) = 0.2$).

\textbf{15.} المعادلتان الناظميتان في السؤال 10 هما
بالضبط $\sum_i\varepsilon_i = 0$ و $\sum_i x_i\varepsilon_i =
0$: أي إن متجهة البواقي \omterm{def:b1:euclid:orthogonal}{متعامدة} مع $C_1 = (1, \dots, 1)$
ومع $C_2 = (x_i)$ --- أي مع فضاء النموذج كله. وعلى
الخصوص يوازن المستقيم الأمثل أخطاءه دائمًا: فمجموعها
معدوم.

\textbf{16.} يعطي $\frac{\dd}{\dd\alpha}\sum(y_i - \alpha)^2 =
-2\sum(y_i - \alpha) = 0$ أن $\alpha = \bar y$ (وتجعله
\omterm{def:b1:derivative:def}{المشتقة} الثانية $2n > 0$ القيمةَ الصغرى الشاملة، لأن
الدالة تربيعية قاسرة في متغير واحد). والقيمة
الصغرى $\sum(y_i - \bar y)^2 = n v_y$: أي إن التباين هو
الخطأ التربيعي غير القابل للاختزال لنموذج ثابت.

\textbf{17.} يعطي $\frac{\dd}{\dd\beta}\sum(y_i - \beta x_i)^2 =
-2\sum x_i(y_i - \beta x_i) = 0$ أن $\beta_0 = \frac{\sum
x_iy_i}{\sum x_i^2}$. وبكتابة الميلين بمقادير
مركَّزة: يساوي $\beta_0 = \frac{c_{xy} + \bar x\bar y}{v_x + \bar
x^2}$ المقدارَ $\frac{c_{xy}}{v_x}$ إذا وفقط إذا كان $c_{xy}\bar x^2 = \bar
x\bar y\,v_x$، أي إذا وفقط إذا كان $\bar x = 0$ (بفواصل مركَّزة) أو
$\bar y = \beta\bar x$ --- والأخير يعني $\alpha = 0$: أي إن
\omterm{pb:b1:multivar:1}{مستقيم الانحدار} الكامل يمرّ أصلًا بالمبدأ.

\textbf{18.} $g(\tau) = \norm{M - P_0}^2 - 2\tau\langle M - P_0,
w\rangle + \tau^2$ أصغري عند $\tau^* = \langle M - P_0,
w\rangle$، بالقيمة $d(M,D)^2 = \norm{M - P_0}^2 - \langle M -
P_0, w\rangle^2$. وفي المستوي، أتمّ $w$ إلى \omterm{def:b1:vspaces:free}{أساس} متعامد
ممنظم $(w, n)$ مع $n = \frac{(a, b)}{\sqrt{a^2+b^2}}$ (وهو الناظم
الواحدي للمقدار $D$): عندئذ $\norm{M - P_0}^2 = \langle M - P_0,
w\rangle^2 + \langle M - P_0, n\rangle^2$، ومنه $d(M, D) =
\abs{\langle M - P_0, n\rangle}$، ومع $aP_{0x} + bP_{0y} =
-c$:
\[
d = \frac{\abs{a x_0 + b y_0 + c}}{\sqrt{a^2 + b^2}} .
\]

\textbf{19.} من $E_{\min}/n = v_y - \frac{c_{xy}^2}{v_x}$ و
$\beta^2 v_x = \frac{c_{xy}^2}{v_x}$:
\[
v_y = \beta^2 v_x + \frac{E_{\min}}n :
\]
أي التباين الكلي $=$ التباين على امتداد المستقيم المقايَس $+$ تباين
البواقي. (وبالقسمة على $v_y$: $1 = \rho^2 + (1 - \rho^2)$، ومنه فحصة
التباين «المفسَّرة» بالمستقيم هي $\rho^2$.)

\textbf{20.} فضاء النموذج هو
$\operatorname{Vect}\bigl(C_1, C_2, C_3\bigr)$ مع $C_1 = (1)$ و
$C_2 = (x_i)$ و $C_3 = (x_i^2)$؛ ويقود تصغير $\norm{aC_1 + bC_2 +
cC_3 - b}^2$، تمامًا كما في السؤال 8، إلى جملة غرام
$3\times3$، ومركّبات مصفوفتها $\langle C_i, C_j\rangle =
\sum_k x_k^{\,i+j-2}$: أي مصفوفة العزوم المعروضة. وهي
قابلة للقلب إذا وفقط إذا كانت $(C_1, C_2, C_3)$ \omterm{def:b1:vspaces:free}{حرة}
(\cref{exo:b1:euclid:11})؛ وعلاقةٌ $a + bx_k + cx_k^2 = 0$
من أجل كل $k$ تجعل كل $x_k$ جذرًا \omterm{def:b1:poly:def}{لكثير حدود} واحد
من درجة $\leq 2$، وهذا محال بثلاث قيم متمايزة ما لم يكن $a
= b = c = 0$. وقارن مصفوفات العزوم في مسألة نهاية الأسبوع
في \cref{ch:b1:det}، حيث كانت محدداتها قيمَ فاندرموند
مربَّعة.

\textbf{21.} الأولى: $r = 2, s = 1, t = 2$، و $rt - s^2 = 3 > 0$ و
$r > 0$: أي قيمة صغرى شاملة قطعية عند المبدأ. والثانية: $r = 2, s
= 3, t = 2$، و $rt - s^2 = -5 < 0$: أي سرج. والثالثة: $r = s = t =
2$، و $rt - s^2 = 0$: أي منحلّة --- لكن $x^2 + 2xy + y^2 = (x +
y)^2 \geq 0$ ينعدم على المستقيم $y = -x$ كله: أي قيمة صغرى
شاملة (غير قطعية) مبلوغة على امتداد مستقيم، غير مرئية
لاختبار المحدد.

\textbf{22.} المجاميع الجديدة ($n = 5$): $\bar x = 3.2$ و $\bar y =
1.6$، و $\sum x_iy_i = 19$ ومنه $c_{xy} = 3.8 - 5.12 = -1.32$،
و $\sum x_i^2 = 114$ ومنه $v_x = 22.8 - 10.24 = 12.56$. والميل الجديد:
\[
\beta = \frac{-1.32}{12.56} \approx -0.105 :
\]
فقد حوّلت نقطة واحدة اتجاهًا متزايدًا بوضوح ($\beta = 1.4$) إلى
اتجاه متناقص قليلًا. ويحمّل الخطأ المربَّع باقيًا
$\varepsilon^2$، ومنه تهيمن نقطة بعيدة واحدة --- بمقدار
$\abs{x_{i} - \bar x}$ كبير \emph{و}باقٍ كبير ---
على $c_{xy}$ و $v_x$ معًا: \omterm{pb:b1:multivar:1}{فالمربعات الصغرى} فعّالة لكنها ليست
متينة.

\textbf{23.} ضع $W = \sum w_i$ وعرّف $\bar x_w =
\frac1W\sum w_ix_i$ و $\bar y_w$ بالمثل و $v_x^w = \frac1W\sum
w_ix_i^2 - \bar x_w^2$ و $c^w_{xy} = \frac1W\sum w_ix_iy_i - \bar
x_w\bar y_w$. \omterm{def:b1:logic:map}{والتطبيق} $\langle u, v\rangle_w = \sum_i w_iu_iv_i$
جداءٌ سلّمي على $\R^n$ (لأن $w_i > 0$ يعطي التعيين)، ومنه
ينطبق الجزء 2 حرفيًا وتنقسم \omterm{pb:b1:multivar:1}{المعادلات الناظمية} على $W$
لتصير
\[
\alpha + \beta\bar x_w = \bar y_w,
\qquad
\alpha\bar x_w + \beta(v^w_x + \bar x_w^2) = c^w_{xy} + \bar
x_w\bar y_w ,
\]
ومنه $\beta = c^w_{xy}/v^w_x$ و $\alpha = \bar y_w -
\beta\bar x_w$: أي الصيغ نفسها، بمتوسطات موزونة كلها.

\textbf{24.} يفرض $E_{\min} = \sum\varepsilon_i^2 = 0$ أن يكون $\varepsilon_i = 0$ من أجل كل دليل: أي أن تقع كل النقاط على المستقيم بالضبط؛ وحسب السؤال
13 هذه هي حالة $\abs\rho = 1$. والاختبار: من أجل $(1,1), (2,3),
(3,5)$: $\bar x = 2$ و $\bar y = 3$ و $v_x = \frac{14}3 - 4 =
\frac23$ و $c_{xy} = \frac{22}3 - 6 = \frac43$: ومنه $\beta = 2$ و
$\alpha = -1$، وفعلًا $y_i = 2x_i - 1$ من أجل النقاط الثلاث؛
و $v_y = \frac{35}3 - 9 = \frac83$ و $\rho^2 =
\frac{(4/3)^2}{(2/3)(8/3)} = 1$.

\textbf{25.} (أ) من أجل الدوال التربيعية يكون النشر من الرتبة الثانية
\emph{متطابقة}، ومنه تصنّف دراسة إشارة الشكل التربيعي
--- وهي جبر محض، السؤالان 1--2 --- النقاطَ الحرجة
دون أيّ مبرهنة تايلور مقبولة. (ب) وتقول \omterm{pb:b1:multivar:1}{المعادلات الناظمية}
«الباقي متعامد مع فضاء النموذج»، ومنه فالقيمة الصغرى
التحليلية \emph{هي} \omterm{thm:b1:euclid:projection}{الإسقاط المتعامد} لمتجهة المعطيات:
فيحسب حساب التفاضل والهندسة الإقليدية الغرضَ نفسه. (ج)
ويقيس الارتباط $\rho = c_{xy}/\sqrt{v_xv_y}$ حصة
$\rho^2$ من التباين التي يفسّرها المستقيم، و
يحدّه كوشي--شوارتز: $\abs\rho \leq 1$، مع المساواة
من أجل المعطيات المستقيمية فقط. (د) وعادت للظهور الأسسُ والحرية (الفصل 18)،
ومصفوفات غرام والعزوم ومحدداتها (الفصلان 22--23)،
\omterm{thm:b1:euclid:projection}{والإسقاط المتعامد} (الفصل 23) بوصفها القطعَ العاملة
في خوارزمية واحدة. ويطوّر الجزآن 2--3
\emph{طريقة \omterm{pb:b1:multivar:1}{المربعات الصغرى}} (لوجاندر وغاوس) عبر
\emph{معادلاتها الناظمية}.
\end{solution}
